% Lösungen Zone (zusammenbau v0.1)
\einheitenkopf[][Kennst du schon]{Das kennst du schon}

Zuordnung: 1 → die Datengrundlage der Prüfungen in Einheit 1 und der fhr-Kette · 2 → der Anteilsvergleich in Einheit 1 und 2 · 3 → die fhr-Prüfungen und die Häufigkeitszeilen · 4 → die Baumaufgaben der Einheit 2 · 5 → die Rückrichtungen der Einheit 2 · 6 → die Vorstellung „ohne Zurücklegen ändert sich p“ der Einheit 3 · 7 → die Datengrundlage der Prüfungen in Einheit 1 und der fhr-Kette · 8 → die Datengrundlage der Prüfungen in Einheit 1 und der fhr-Kette

\erg{1}{a) 6, 4, 10, 18, 12, 30, 24, 16, 40 (Zeilen: Brille, keine Brille, Summe; Spalten: Mädchen, Jungen, Summe) \quad b) 18, 12, 30, 37, 33, 70, 55, 45, 100 (Prozent; Zeilen: Rad, kein Rad, Summe; Spalten: weiblich, männlich, Summe)}

1a)

\vierfeldertafel{Mädchen}{Brille}{6,4,10,18,12,30,24,16,40}


1b)

\vierfeldertafel{Weiblich}{Rad}{18\,\%,12\,\%,30\,\%,37\,\%,33\,\%,70\,\%,55\,\%,45\,\%,100\,\%}

\erg{2}{a) $\frac{5}{8} = 0{,}625$ \quad b) $P_A(B) = \frac{0{,}1}{0{,}4} = 0{,}25$}
\erg{3}{a) $\frac{12}{48} = \frac{1}{4} = 0{,}25$ \quad b) $\frac{875}{1\,250} = 0{,}7$}
\erg{4}{a) $p^2$ \quad b) $2 \cdot p \cdot (1 - p)$ – die Pfade Rot-Blau und Blau-Rot addieren}
\erg{5}{a) $x = 0{,}4$ \quad b) $x^2 - 0{,}5x + 0{,}06 = 0$, also $x_1 = 0{,}2$ und $x_2 = 0{,}3$}
\erg{6}{a) $\frac{1}{6} \cdot \frac{1}{6} = \frac{1}{36}$ \quad b) $\frac{2}{5} \cdot \frac{1}{4} = \frac{1}{10}$}
\erg{7}{a) Fehler: Die 30\,\% beziehen sich auf die 20 Jugendlichen, nicht auf alle 50 Mitglieder. Richtig: $0{,}3 \cdot 20 = 6$ Jugendliche spielen Tennis, 9 Erwachsene.}
\erg{8}{a) $0{,}25 \cdot 32 = 8$ Kinder, also $30 - 8 = 22$ Erwachsene}
